% !TEX root = ../presentation.tex
% =============================================================
%  Subsection — Steady-state angular stability of an AC OHL
%  Part of ac-transmission; \input from ac-transmission.tex,
%  between transmission-limits.tex and series-compensation.tex.
%
%  Figures in ./fig (as_*.pdf), all vector, generated by
%  fig/angular_stability_figures.py (numpy + matplotlib, no scipy)
%  from the same exact two-port model as the other sub-files:
%    as_power_angle.pdf        P(delta) and K_s of the 300 mi line
%    as_electrical_length.pdf  P_max vs length: exact / lossless / lumped
%    as_contingency.pdf        two parallel circuits, intact and after N-1
%    as_equal_area.pdf         equal-area criterion (generic machine p.u.)
%  The same script prints every number quoted on the slides.
%
%  Reference line: 500 kV, 60 Hz, 3x Bluebird ACSR, flat, 35 ft / 18 in,
%  Z_w = 266 Ohm, P_nat = 939 MW, X_L = 162.4 Ohm, beta*l = 35.2 deg.
% =============================================================
\providecommand{\acfig}[1]{ac-transmission/fig/#1}

\intermezzo{Angular Stability}{Why the transmission angle is held far below $90^\circ$}

% ---------------------------------------------------------------
\begin{frame}{The power--angle relation and the synchronising power}
  \begin{columns}[c,onlytextwidth]
    \begin{column}{0.53\textwidth}\centering
      \topoimg[width=\linewidth,height=0.58\textheight,keepaspectratio]{\acfig{as_power_angle.pdf}}
    \end{column}
    \begin{column}{0.45\textwidth}
      \setlength{\abovedisplayskip}{2pt}\setlength{\belowdisplayskip}{2pt}
      \[ P = \frac{U_S\cdot U_R}{X}\cdot\sin\delta \]
      \[ K_s = \frac{\partial P}{\partial\delta}
             = \frac{U_S\cdot U_R}{X}\cdot\cos\delta \]
      {\footnotesize
      \begin{itemize}
        \setlength{\itemsep}{0.35em}
        \item $K_s$ is the stiffness of the coupling. Below $P_{max}$ a
              machine that advances delivers more power and is pulled
              back; beyond it, less --- and it pulls out.
        \item $K_s=0$ at $P_{max}$ is the steady-state stability limit.
              With losses it sits at $81^\circ$, not $90^\circ$.
        \item At $\delta=30^\circ$: $0.81\cdot P_{nat}$, at $79\,\%$ of
              the stiffness at $\delta=0$.
      \end{itemize}}
    \end{column}
  \end{columns}
  \atnote{Exact two-port, both terminals at 1.0\,p.u., 300\,mi reference
  line; the lossless forms above are shown for orientation.}
\end{frame}

% ---------------------------------------------------------------
\begin{frame}{Electrical length, not lumped reactance}
  \begin{columns}[c,onlytextwidth]
    \begin{column}{0.55\textwidth}\centering
      \topoimg[width=\linewidth,height=0.60\textheight,keepaspectratio]{\acfig{as_electrical_length.pdf}}
    \end{column}
    \begin{column}{0.43\textwidth}
      \setlength{\abovedisplayskip}{3pt}\setlength{\belowdisplayskip}{5pt}
      \[ X' = Z_w\cdot\sin\left(\beta\cdot l\right) \]
      \[ P_{max} = \frac{U^2}{X'}
                 = \frac{P_{nat}}{\sin\left(\beta\cdot l\right)} \]
      {\footnotesize
      \begin{itemize}
        \setlength{\itemsep}{0.35em}
        \item $\beta=\omega\cdot\sqrt{L'\cdot C'}$: at 60\,Hz the
              wavelength is 4937\,km, a quarter wave 767\,mi (1234\,km).
        \item At a quarter wave $P_{max}=P_{nat}$ exactly --- the
              analytical spine under the St.~Clair curve.
        \item 300\,mi: lumped $\omega\cdot L'\cdot l=162.4\,\Omega$
              against $X'=153.4\,\Omega$, and $1/\sin\left(\beta\cdot l\right)
              =1.73$ against 1.45 exact.
      \end{itemize}}
    \end{column}
  \end{columns}
  \atnote{Lumped $X$ overstates beyond roughly 250\,km; $R=27.5\,\Omega$
  and the line charging account for the rest of the gap at 300\,mi.}
\end{frame}

% ---------------------------------------------------------------
\begin{frame}{The reactance is not only the line}
  \begin{columns}[T,onlytextwidth]
    \begin{column}{0.50\textwidth}
      {\small
      \begin{itemize}
        \setlength{\itemsep}{0.4em}
        \item The $X$ of the power--angle relation is the whole path
              between the two internal EMFs, not the line alone.
        \item A 300\,mi line into systems of 10\,GVA at each end loses a
              quarter of its $P_{max}$ to the terminals.
        \item Loadability quoted for a line in isolation is therefore
              optimistic; the short-circuit level at both ends belongs in
              the statement.
        \item The same quantity drives converter performance, where it is
              normalised as the SCR.
      \end{itemize}}
    \end{column}
    \begin{column}{0.47\textwidth}
      \setlength{\abovedisplayskip}{2pt}\setlength{\belowdisplayskip}{4pt}
      \[ X_{tot} = X_{th\_s} + X_{tr\_s} + X' + X_{tr\_r} + X_{th\_r} \]
      \[ \mathrm{SCR} = \frac{S_{sc}}{P_n},\qquad
         X_{th} = \frac{U_n^2}{S_{sc}} \]
      \renewcommand{\arraystretch}{1.15}
      {\footnotesize
      \begin{tabular}{@{}l r r r@{}}
        \textbf{SCL, each end} & \multicolumn{1}{c}{\textbf{$X_{th}$}}
                               & \multicolumn{1}{c}{\textbf{$X_{tot}$}}
                               & \multicolumn{1}{c}{\textbf{$P_{max}$}} \\
        \multicolumn{1}{c}{GVA} & \multicolumn{1}{c}{$\Omega$}
                               & \multicolumn{1}{c}{$\Omega$}
                               & \multicolumn{1}{c}{MW} \\
        \hline
        5        & 50.0 & 253.4 &  986 \\
        10       & 25.0 & 203.4 & 1229 \\
        20       & 12.5 & 178.4 & 1401 \\
        stiff    &  0.0 & 153.4 & 1630 \\
        \hline
      \end{tabular}}
    \end{column}
  \end{columns}
  \atnote{Lossless $U^2/X_{tot}$, transformers neglected, 300\,mi
  reference line with $X'=153.4\,\Omega$. $S_{sc}$ short-circuit level,
  $P_n$ rated power of the infeed.}
\end{frame}

% ---------------------------------------------------------------
\begin{frame}{From the maximum to the operating limit}
  \begin{columns}[c,onlytextwidth]
    \begin{column}{0.55\textwidth}\centering
      \topoimg[width=\linewidth,height=0.60\textheight,keepaspectratio]{\acfig{as_contingency.pdf}}
    \end{column}
    \begin{column}{0.43\textwidth}
      {\footnotesize
      \begin{itemize}
        \setlength{\itemsep}{0.4em}
        \item Two 300\,mi circuits: 2726\,MW intact, 1363\,MW on the
              survivor after N$-$1. The reactance of the path doubles.
        \item Planning criterion
              $\left(P_{max}-P\right)/P_{max}\ge0.30$ on the
              \emph{post-contingency} network $\rightarrow$ 954\,MW, i.e.
              $1.02\cdot P_{nat}$ for the pair.
        \item 1200\,MW looks comfortable intact at $\delta=22.9^\circ$
              and sits $12\,\%$ below the post-contingency maximum ---
              not acceptable.
        \item The thermal rating of the pair is 6755\,MW: seven times the
              binding constraint.
      \end{itemize}}
    \end{column}
  \end{columns}
  \atnote{Both terminals at 1.0\,p.u., uncompensated, exact two-port.
  Series compensation raises the post-contingency curve --- see the next
  section.}
\end{frame}

% ---------------------------------------------------------------
\begin{frame}{Transient stability --- the equal-area criterion}
  \begin{columns}[c,onlytextwidth]
    \begin{column}{0.55\textwidth}\centering
      \topoimg[width=\linewidth,height=0.54\textheight,keepaspectratio]{\acfig{as_equal_area.pdf}}
    \end{column}
    \begin{column}{0.43\textwidth}
      \setlength{\abovedisplayskip}{2pt}\setlength{\belowdisplayskip}{2pt}
      \[ \begin{aligned}
           A_{acc} &= \int_{\delta_0}^{\delta_{cc}}\left(P_m-P_{flt}\right)\mathrm{d}\delta \\
           A_{dec} &= \int_{\delta_{cc}}^{\delta_{max}}\left(P_{pst}-P_m\right)\mathrm{d}\delta
         \end{aligned} \]
      {\footnotesize
      \begin{itemize}
        \setlength{\itemsep}{0.3em}
        \item The output collapses during the fault while the mechanical
              input continues; the rotor must repay that energy before
              $\delta_{max}$.
        \item A higher pre-fault $\delta$ shrinks $A_{dec}$ from both
              ends, so the critical clearing time falls steeply with
              loading.
        \item Fast clearing --- two-cycle breakers, independent-pole
              tripping --- buys capability.
      \end{itemize}}
    \end{column}
  \end{columns}
  \atnote{Generic machine p.u.: $P_m=0.6$, $P_{max}$ 1.0 pre-fault, 0.3
  during the fault, 0.7 after it. $\delta_0=36.9^\circ$,
  $\delta_{cc}=45.5^\circ$, $\delta_{max}=121.0^\circ$.}
\end{frame}

% ---------------------------------------------------------------
\begin{frame}{Damping sets the limit before $P_{max}$ does}
  \begin{columns}[T,onlytextwidth]
    \begin{column}{0.50\textwidth}
      {\small
      \begin{itemize}
        \setlength{\itemsep}{0.4em}
        \item As $\delta$ grows $K_s$ falls: the electromechanical mode
              slows and its damping ratio drops.
        \item High-gain AVRs --- what allows a large $\delta$ in the
              first place --- give negative damping torque at heavy
              loading. The PSS exists to put it back.
        \item In a multi-machine system the limit is the eigenvalue
              criterion, not $P_{max}$: it can bind well below the
              two-machine maximum.
      \end{itemize}}
    \end{column}
    \begin{column}{0.47\textwidth}
      \setlength{\abovedisplayskip}{2pt}\setlength{\belowdisplayskip}{4pt}
      \[ \frac{2\cdot H}{\omega_0}\cdot\Delta\ddot{\delta}
         + K_D\cdot\Delta\dot{\delta} + K_s\cdot\Delta\delta = 0 \]
      \[ \omega_n = \sqrt{\frac{K_s\cdot\omega_0}{2\cdot H}} \]
      \renewcommand{\arraystretch}{1.15}
      {\footnotesize
      \begin{tabular}{@{}l r@{}}
        \textbf{Mode} & \multicolumn{1}{c}{\textbf{Hz}} \\
        \hline
        Inter-area           & 0.2--0.8 \\
        Local plant          & 0.8--2.0 \\
        Torsional, SSR/SSCI  & 10--50 \\
        \hline
        \multicolumn{2}{@{}l}{Criterion: damping ratio $\ge 3\ldots5\,\%$}\\
      \end{tabular}}
    \end{column}
  \end{columns}
  \atnote{$H$ inertia constant, $K_D$ damping torque coefficient,
  $\omega_0$ rated angular frequency. Inter-area modes are the ones that
  constrain long transfers.}
\end{frame}

% ---------------------------------------------------------------
\begin{frame}{A DC link has no transmission angle}
  \begin{columns}[T,onlytextwidth]
    \begin{column}{0.48\textwidth}
      {\small
      \begin{itemize}
        \setlength{\itemsep}{0.4em}
        \item No angle, no synchronising torque, no equal-area criterion:
              the converter control sets power as a commanded quantity.
        \item Length enters only through the conductor resistance, so the
              $1/\sin\left(\beta\cdot l\right)$ decay has no counterpart.
        \item An embedded HVDC link can be modulated to \emph{damp} the
              inter-area modes of the surrounding AC system instead of
              loading them.
        \item The constraint moves to the terminals: SCR and commutation
              robustness for LCC, control interaction for VSC.
      \end{itemize}}
    \end{column}
    \begin{column}{0.49\textwidth}
      \renewcommand{\arraystretch}{1.25}
      {\footnotesize
      \begin{tabularx}{\linewidth}{@{}l
          >{\raggedright\arraybackslash}X
          >{\raggedright\arraybackslash}X@{}}
        & \textbf{AC line} & \textbf{HVDC link} \\
        \hline
        Power set by  & $\delta$ across $X$ & converter control \\
        Length        & $1/\sin\beta l$     & $R$ only \\
        Synchronising & inherent            & by control \\
        Binding limit & stability           & AC strength, converter
                                              rating \\
        \hline
      \end{tabularx}}
    \end{column}
  \end{columns}
  \atnote{The angular limits of the preceding slides are the reason a
  long AC corridor needs compensation and a DC corridor does not.}
\end{frame}
